% Transcription — fonds Alexandre Grothendieck, shelfmark 10, batch 6 (pages 101-120).
% Established from the facsimile archives/batches/10/batch-06.pdf (a mirror of
% https://grothendieck.umontpellier.fr/10.pdf), per the skill
% transcribe-grothendieck.
%
% Pass: Opus 5 (claude-opus-5), 2026-09-12 — first pass, unchecked against the
% pages by a human.
%
% Every leaf of this batch is on the subject of the carton, and every one of
% them is typed — three typescripts, all corrected in his hand — except that
% the middle of the lot, pages 110 to 112, is three leaves of his own.
%
% Pages 101-108: the typescript announced by the « Notes Saavedra » of page 99,
% headed « Quelques exemples de catégories tensorielles », in six numbered
% items. It is addressed to Saavedra — page 107 says « Je laisse le soin à
% Saavedra de déterminer … » — and it is the only run of this folder that
% names him.
%
%   101  1) the category C_M of M-graded finite-dimensional vector spaces, its
%        forgetful fibre functor and its group D_k(M); the application, that a
%        functorial M-grading of a fibre functor F is a homomorphism
%        D_k(M) → G; and the criterion for such a grading to come from a
%        grading of the objects themselves — that the homomorphism be central,
%        with the canonical central i: G_m → G on motives, the Tate motive
%        giving j: G → G_m = Gl(1) and the weight relation ji(λ) = λ².
%   102  the Hodge fibre functor in characteristic zero, G_m² → G, i.e. two
%        commuting i_1, i_2 with i(λ) = i_1(λ)i_2(λ), and the warning that
%        i_1, i_2 are not central; 2) the De Rham functor, graded and filtered
%        rather than bigraded, the degeneration of the hypercohomology spectral
%        sequence giving Gr(H_DR(V)) ≃ H_Hdg(V).
%   103  filtrations of a fibre functor, decreasing, discrete, indexed by Z and
%        compatible with ⊗; why such a datum is not a homomorphism into G (the
%        category of filtered vector spaces is not abelian); the second fibre
%        functor F' = Gr F, the torsor P under G, the grading i: G_m → G', the
%        subgroup H'_{i_1} of filtration-respecting automorphisms inducing the
%        identity on the graded object, and Q = Isom^filt(F,F') as a pseudo-
%        torsor under H'.
%   104  the resulting dictionary: a filtered fibre functor = a fibre functor
%        F', plus a grading i_1: G_m → G', plus a right torsor Q' under
%        H'_{i_1}, with F = F' ∧^{H'} Q'; then, cancelled by diagonals, the NB
%        on the De Rham functor on motives; and, uncancelled, the statement
%        that H' is always a projective limit of unipotent groups, so that over
%        a countably generated category every torsor under it is trivial and
%        every filtered fibre functor is in fact a graded one.
%   105  the splitting as a point of a torsor under the pro-unipotent H, « pas
%        du tout canonique »; 3) pré-structures de Hodge on a Q-vector space,
%        the two conditions a) and b), the category they form and its Galois
%        group; the bigrading as i_1, i_2: G_{m,C} → G_C, condition a) saying
%        that i_C = i_1i_2 is defined over Q and central.
%   106  condition b) saying i_2 = ī_1; the Betti-Hodge functor from motives
%        over C to pré-structures, the standard Hodge conjecture read as its
%        full faithfulness, i.e. as the surjectivity of G → G_0; polarisation
%        of weight n as a pairing φ: V × V → Q(n) whose associated hermitian
%        form ψ(x,y) = φ(x,ȳ)(-i)^{p-q} is positive definite, his hand
%        correcting the bidegree to (p, n-p).
%   107  the image of the Betti-Hodge functor, « on n'a aucune idée » of what
%        it is; pré-structures being semi-simple, the exercise of writing the
%        Hodge condition on the Galois group « plaisant et délectable »;
%        4) what Saavedra is left to determine — the real structure of Hodge
%        deduced from X_R = X_0, and the Frobenius at infinity f_∞ of order 2.
%   108  5) the criterion for the Galois group to be profinite — every object
%        generating a semi-simple subcategory with finitely many simple objects
%        — and C_0 as the subcategory answering the pro-group of the G_i°;
%        6) polarisation on a general ⊗-category over a subfield of R, left for
%        another occasion.
%
% Pages 110-112, his own hand: three leaves on the categories of modules with
% operators of an affine group scheme. Page 110 fixes the hypotheses on G and
% the equivalence C(G) ≃ Ind C_f(G) (SGAD VI_B 11.2), then asks for the
% category of G_T-modules in terms of C(G) and B = f_*(O_T); page 111 answers
% it through the Weil restriction of Aut(F); page 112 turns the question round
% — to what extent is G recovered from (C_f(G), φ_G^f) as a ⊗-category with a
% ⊗-functor — and closes on the question, marked « ??? », of which ⊗-categories
% come from an affine group scheme.
%
% Page 114: the plan of the book in nine chapters, with an appendix and two
% further chapters added in his hand — one on the non-commutative case, one on
% Picard categories and on categories carrying both ⊕ and ⊗.
%
% Pages 115-116: the detailed table of contents of Chapter 1, annotated
% throughout — the paragraph numbers renumbered by hand, two subsections
% added, one struck, a « non ! » against the identification of Hom with Moh,
% and a margin list of five examples of ⊗-laws.
%
% Pages 117-120: a fourth typescript, headed « ⊗-catégories », giving the
% axioms — associativity, commutativity, unit, internal Hom, biadditivity —
% then ⊗-functors, homomorphisms of ⊗-functors, the dissociation of the
% datum of ⊗ into independent aspects, the monoid End(1), and the
% bialgebra/coalgebra dictionary of item 6).
%
% Pages 109 and 113 carry nothing: they are absent from what follows.
%
% Two conventions of these typescripts are recorded once and not repeated. The
% machine had no ⊗ and no Greek: ⊗ is struck over an O, φ over an o, ψ over a
% y, and the primes are apostrophes which the carbon thickens. And « gauche »
% is overtyped by « droite » at every occurrence in pages 103-105 — the
% typist's correction, not his — so « droite » is what is given.
%
% The dating is the inventory's own, for the whole shelfmark.
\documentclass[11pt,a4paper]{article}
% Le préambule commun à toutes les transcriptions du fonds Grothendieck.
%
% Il tient en une idée : **l'incertitude se compose, elle ne se résout pas.**
% Une transcription qui lisse un mot douteux a détruit la seule chose pour
% laquelle on la faisait — un lecteur doit pouvoir savoir, à chaque endroit,
% ce qui était sur le feuillet et ce qui a été ajouté. Les six macros qui
% suivent sont donc l'appareil critique, et non de la décoration.
%
% Les mêmes commandes sont reconnues par `scripts/render.mjs`, qui produit la
% vue de lecture du site. Une macro ajoutée ici doit l'être aussi là-bas,
% faute de quoi le rendu échouera bruyamment — ce qui est voulu : un rendu
% silencieusement faux est pire qu'un rendu absent.

\ProvidesPackage{grothendieck}[2026/08/08 Transcriptions du fonds Grothendieck]

\usepackage{amsmath,amssymb,amsthm}
\usepackage{mathrsfs}
\usepackage{stmaryrd}
% Les diagrammes commutatifs sont partout dans ce fonds : c'est la langue
% même de la géométrie algébrique de Grothendieck, et une transcription qui
% les rendrait en prose ne transcrirait rien.
\usepackage{tikz-cd}
% Français par défaut — c'est la langue des feuillets — mais l'anglais est
% chargé aussi : la traduction et le résumé sont des documents anglais, et la
% typographie française (espaces avant les deux-points) y serait une faute.
% Ces documents basculent par \selectlanguage{english} en tête de corps.
\usepackage[english,french]{babel}
\usepackage{fontspec}
\usepackage{geometry}
\geometry{a4paper,margin=25mm,marginparwidth=28mm}
\usepackage{marginnote}
\usepackage{xcolor}
% ulem, not soul. `soul`'s \st and \ul are built on a character-by-character
% scan that XeLaTeX's font machinery breaks: the document fails with an
% undefined control sequence rather than degrading. ulem does the same two jobs
% and is Unicode-engine safe. `normalem` stops it hijacking \emph, which the
% transcriptions use for Grothendieck's own emphasis.
\usepackage[normalem]{ulem}
\usepackage{hyperref}
\hypersetup{colorlinks=true,linkcolor=black,urlcolor=black,pdfborder={0 0 0}}

% --- Métadonnées du lot -----------------------------------------------------
% Reprises telles quelles par le script de rendu, qui les affiche en tête de
% la vue de lecture. Elles fixent ce que le fichier prétend couvrir : sans
% elles, une transcription flotte sans point d'attache dans un fonds de seize
% mille feuillets.

\newcommand{\folder}[1]{\gdef\grothFolder{#1}}
\newcommand{\batch}[1]{\gdef\grothBatch{#1}}
\newcommand{\leaves}[2]{\gdef\grothFirst{#1}\gdef\grothLast{#2}}
\newcommand{\foldertitle}[1]{\gdef\grothFoldertitle{#1}}
\gdef\grothFolder{?}\gdef\grothBatch{?}\gdef\grothFirst{?}\gdef\grothLast{?}\gdef\grothFoldertitle{}

% --- Le résumé --------------------------------------------------------------
%
% Il ouvre la lecture modernisée, sous le titre « Résumé », et il est composé
% en petit. Ce n'est pas un ornement : c'est le seul passage du document qui
% ne soit pas des mathématiques, et un lecteur doit pouvoir le reconnaître
% avant de l'avoir lu — puis le sauter s'il connaît déjà le sujet, ou s'y
% arrêter s'il ne le connaît pas. Le filet qui le suit marque la reprise du
% corps.

\newenvironment{resume}
  {\small\setlength{\parskip}{0.45em}}
  {\par\medskip\hrule\medskip}

% --- Mots-clés ---------------------------------------------------------------
%
% En anglais, en clôture du résumé de la lecture modernisée : le vocabulaire
% moderne sous lequel chercher ce que les feuillets construisent. Le manifeste
% du site (`npm run manifest`) les extrait pour étiqueter la cote entière —
% cette ligne est la source unique des tags, il n'y a pas de fichier à côté.
\newcommand{\keywords}[1]{\par\smallskip\noindent
  {\footnotesize\textsc{Keywords} --- \textit{#1}}\par}

% --- L'appareil critique ----------------------------------------------------

% Le numéro de feuillet, en marge. C'est l'ancre qui permet de régler un
% désaccord sur une formule en regardant *un* feuillet plutôt qu'une chemise
% de six cents. Le site s'en sert aussi pour tourner les pages du fac-similé
% quand on fait défiler la transcription.
\newcommand{\leaf}[1]{%
  \marginnote{\footnotesize\color{gray!70}\textsf{#1}}%
  \label{leaf:#1}\ignorespaces}

% Illisible : on ne devine pas. Un mot inventé qui se lit comme les autres est
% le pire résultat possible.
\newcommand{\ill}{\textcolor{red!60!black}{[\,\ldots\,]}}

% Lecture douteuse : le mot est donné, et signalé comme incertain.
\newcommand{\uncertain}[1]{\uline{#1}}

% Ajout de l'éditeur — une lettre manquante, un mot sous-entendu.
\newcommand{\add}[1]{\textcolor{blue!50!black}{[#1]}}

% Biffé par Grothendieck lui-même. Ce qu'il a rayé fait partie du document :
% une rature dit souvent où la pensée a changé de direction.
\newcommand{\struck}[1]{\sout{#1}}

% Note du transcripteur, sur l'état matériel du feuillet.
\newcommand{\note}[1]{\par\smallskip{\small\color{gray!60!black}\textit{[#1]}}\par\smallskip}

% Note marginale de Grothendieck, distincte du corps du texte.
\newcommand{\marginal}[1]{\par\smallskip\noindent\hspace{1em}%
  {\small\color{gray!60!black}#1}\par\smallskip}

% --- Titre ------------------------------------------------------------------

\AtBeginDocument{%
  \begin{center}
    {\large\bfseries Fonds Alexandre Grothendieck --- cote n\textsuperscript{o}~\grothFolder}\\[2pt]
    {\itshape\grothFoldertitle}\\[4pt]
    {\small feuillets \grothFirst--\grothLast\ (lot \grothBatch)}\\[2pt]
    {\footnotesize\color{gray!60!black}Transcription établie d'après le fac-similé numérisé
     par l'Université de Montpellier.\\
     Les lectures incertaines sont soulignées, les illisibles notées [\,\ldots\,],
     les ajouts entre crochets.}
  \end{center}
  \vspace{1em}}

\endinput


\folder{10}
\batch{6}
\pages{101}{120}
\dating{[à partir de 1958]}
\watermark{Édition de démonstration}
\foldertitle{Catégories tensorielles : notes manuscrites (s.d.), tapuscrits (s.d.)}

\begin{document}

\section*{Quelques exemples de catégories tensorielles (pages 101 à 108)}

\note{Tapuscrit annoncé par les deux mots « Notes Saavedra » que porte, de sa
main, la page 99. Il est corrigé de bout en bout à l'encre : ses insertions
sont données entre crochets, ses suppressions barrées, ses notes marginales en
retrait. La machine n'avait ni $\otimes$ ni lettres grecques ; on ne relève
pas un à un les signes encrés de sa main.}

\page{101}

\subsection*{Quelques exemples de catégories tensorielles.}

1) Soit $M$ un groupe. Soit $\underline{C}_{M}$ la catégorie des vectoriels de
dimension finie sur le corps $k$, munis d'une graduation de type $M$. C'est
une catégorie tensorielle sur $k$, munie d'un foncteur fibre sur $k$, le
foncteur « oubli de la graduation ». Le groupe algébrique associé est le
groupe de type multiplicatif $D_{k}(M)$ (SGA 3 I 4.7.3). Par exemple si
$M = \underline{\mathbb{Z}}^{r}$, on trouve $G = \underline{G}_{m}^{r}$.

Application : soit $\underline{C}$ une catégorie tensorielle sur $k$ munie
d'un foncteur fibre $F$ sur $k$, donc associée à un schéma en groupes affine
$G$ sur $k$. On cherche toutes les façons de mettre, pour chaque
$M \in \operatorname{Ob} \underline{C}$, une graduation de type $M$ sur
$F(V)$, de façon fonctorielle en $M$, et compatible (dans un sens évident)
avec les produits tensoriels.
\note{La machine porte bien « pour chaque $M \in \operatorname{Ob}
\underline{C}$ » puis « fonctorielle en $M$ », là où l'objet variable est $V$
et où $M$ est le groupe des degrés ; on ne corrige pas.}
Elles correspondent aux $\otimes$-foncteurs de $\underline{C}$ dans
$\underline{C}_{M}$ compatibles avec les foncteurs fibres, donc aux
homomorphismes de $D_{k}(M)$ dans $G$. Par exemple, si
$M = \underline{\mathbb{Z}}^{r}$, il faut prendre les homomorphismes de
$\underline{G}_{m}^{r}$ dans $G$.

Dans la situation précédente, on peut se demander quand une graduation de type
$M$ du foncteur $F$ correspond à une graduation de type $M$ du foncteur
identique de $\underline{C}$, i.e.\ pour tout $V$, la graduation de $F(V)$
provient d'une graduation de $V$. On trouve qu'il faut et il suffit pour cela
que l'homomorphisme correspondant $D_{k}(M) \to G$ soit central. Par exemple,
si $\underline{C}$ est la catégorie des motifs sur $k$, et si on dispose d'un
foncteur fibre $F$ de $\underline{C}$ sur $k$, alors on trouve un
homomorphisme central canonique $i : \underline{G}_{m} \to G$. D'ailleurs, la
donnée du motif de Tate (qui est de rang $2$, et de « poids » $2$ pour la
graduation naturelle) correspond à la donnée d'un homomorphisme
$j : G \to \underline{G}_{m} = Gl(1)$.
\note{Le tapuscrit porte « de rang $2$ » ; le motif de Tate est de rang $1$.
Rien sur la page ne corrige.}
Le fait que $T$ soit de poids $2$ s'exprime par la relation
\[
j\,i(\lambda) = \lambda^{2} .
\]

\page{102}
Lorsque $k$ est de caractéristique nulle, on a toujours un foncteur fibre
naturel : le \emph{foncteur de Hodge}, qui à la cohomologie motivique d'une
variété (projective lisse) $X$ associe le vectoriel bigradué
$\coprod H^{q}(X, \underline{\Omega}^{p}_{X/k})$. Donc pour le groupe de
Galois motivique correspondant $G$, on trouve un homomorphisme naturel
\[
\underline{G}_{m}^{2} \longrightarrow G ,
\]
i.e.\ deux homomorphismes commutant l'un à l'autre
\[
i_{1}, i_{2} : \underline{G}_{m} \longrightarrow G .
\]
Le fait que la graduation totale dans la cohomologie de Hodge corresponde au
poids des motifs s'exprime par la relation
\[
i(\lambda) = i_{1}(\lambda)\, i_{2}(\lambda) \, ;
\]
on fera attention que $i_{1}$ et $i_{2}$ ne sont pas centraux (car la
bigraduation en cohomologie de Hodge ne correspond pas à une bigraduation d'un
motif !)

2) Prenant toujours pour $\underline{C}$ la catégorie des motifs sur $k$, avec
$k$ de car.\ nulle, on a un autre foncteur fibre canonique, le \emph{foncteur
de De Rham}, qui associe à la cohomologie motivique d'une variété $X$ le
vectoriel $H^{*}(X, \underline{\Omega}^{*}_{X/k})$ (espace
d'hypercohomologie). Ce vectoriel n'est plus bigradué, mais seulement gradué
et filtré, la filtration étant celle associée à la suite spectrale
d'hypercohomologie, commençant avec la cohomologie de Hodge. On sait
d'ailleurs que cette suite spectrale dégénère (théorie de Hodge), donc
$\operatorname{Gr}(H_{DR}(V)) \simeq H_{Hdg}(V)$
\marginal{où on gradue le 2ème membre par $p$}
(isomorphisme fonctoriel en le motif $V$). On peut se proposer d'analyser à
quelle structure supplémentaire, sur le groupe de Galois motivique associé au
foncteur fibre $H_{DR}$, correspond la filtration canonique de ce foncteur.

De façon générale, étant donné une catégorie tensorielle $\underline{C}$ sur
$k$ munie d'un foncteur fibre $F$, on peut se proposer de déterminer les
filtra\-

\page{103}
tions sur $F$ (décroissantes, \add{discrètes,} indexées par
$\underline{\mathbb{Z}}$) compatibles avec les produits tensoriels (en
utilisant la notion évidente de produit tensoriel de deux espaces vectoriels
filtrés). On notera qu'une telle donnée ne pourra plus s'exprimer par un
homomorphisme d'un certain groupe algébrique dans $G$ (le groupe de Galois de
$\underline{C}$ en $F$), car la catégorie des espaces vectoriels de dimension
finie sur $k$, munis d'une filtration décroissante discrète indexée par
$\underline{\mathbb{Z}}$, n'est pas une catégorie abélienne (les bimorphismes
ne sont pas des isomorphismes). Mais à une telle donnée est associé un
deuxième foncteur fibre \add{$F'(V) =$} $\operatorname{Gr} F(V)$, à valeurs
cette fois-ci dans les vectoriels gradués (N.B.\ $F$ joue le rôle de la
cohomologie de \struck{Hodge} De Rham, $F'$ celui de la cohomologie de Hodge,
muni de la graduation \struck{totale} \add{par $p$}). La donnée de $F'$
correspond alors à un torseur \add{(à droite)} $P$ sous $G$, le groupe de
Galois correspondant $G'$ étant déduit de $G$ en tordant par $P$. La
graduation de $F'$ correspond à un homomorphisme
$i : \underline{G}_{m} \to G'$. Considé\-rons alors sur $F'$ la filtration
décroissante associée à sa graduation, et soit
\[
H'_{i_{1}} = \underline{\operatorname{Aut}}\ \underline{\operatorname{filt}}^{1}(F')
\subset \underline{\operatorname{Aut}}(F') = G'
\]
\note{Le signe porté en exposant après « filt » est un trait vertical
surmontant un point ; il est lu ici comme un $1$, sans que la page le décide.
Le même signe revient à « Isom filt » deux lignes plus bas.}
le sous-schéma en groupes de $G$ qui correspond aux \add{$\otimes$-}auto\-mor\-phismes
de $F'$ (ou plutôt des $F'_{k'}$, $k'$ une $k$-algèbre quelconque) qui
respectent sa filtration \add{(et induisant l'identité sur le gradué
associé)} ; il est canoniquement déterminé par $i_{1}$. On peut aussi regarder
le sous-schéma
\[
Q = \underline{\operatorname{Isom}}\ \underline{\operatorname{filt}}^{1}(F,F')
\subset P = \underline{\operatorname{Isom}}(F,F')
\]
du schéma $P$ des $\otimes$-isomorphismes de $F$ avec $F'$, \struck{formé de}
qui correspond aux automorphismes respectant les filtrations de $F$ et $F'$
\add{(et induisant l'identité sur les gradués associés)}. C'est à priori un
pseudo-torseur à droite sous $H'$, i.e.\ il est vide ou un torseur à droite
sous $H'$.
\note{« droite » est ici frappé par-dessus « gauche », et de même à chacune de
ses occurrences jusqu'à la page 105 ; on donne le mot corrigé.}
Il faudrait prouver que c'est bien un torseur (i.e.\ que sur une extension
convenable $k'$ de $k$, on peut trouver un isomorphisme de $F_{k'}$ avec
$F'_{k'}$ respectant les filtrations). Donc on trouve une restriction du
groupe d'opérateurs (à droite) $G'$ de $P$ au sous-groupe $H'$,

\page{104}
par un $H'$-torseur à droite $Q$.

Moyennant la vérification laissée en suspens à l'instant, on trouve alors que
la donnée d'un « foncteur fibre \emph{filtré} » $F$ de $\underline{C}$ sur $k$
revient à la donnée

\begin{enumerate}
\item[(i)] d'un foncteur fibre $F'$ \struck{gradué de type
$\underline{\mathbb{Z}}$} (d'où un groupe de Galois
$G' = \underline{\operatorname{Aut}}(F')$) ;
\item[(ii)] d'une graduation de type $\underline{\mathbb{Z}}$ de $F'$, i.e.\ un
homomorphisme
\[
i_{1} : \underline{G}_{m} \longrightarrow G' ,
\]
ce qui permet de définir comme ci-dessus un sous-groupe $H'_{i_{1}}$ de $G'$ ;
\item[(iii)] d'un torseur à droite $Q'$ sous $H'_{i_{1}}$.
\end{enumerate}

A ces données, on associe simplement le foncteur fibre tordu
\[
F = F' \overset{H'_{i_{1}}}{\wedge} Q' ,
\]
$F$ étant filtré par la filtration déduite de celle de $F'$
\struck{\ill{}} (associée à la graduation de $F'$) en tordant par $Q'$.

\struck{N.B. Dans le cas où $F$ est le foncteur de De Rham sur la catégorie
des motifs sur $k$, on prouve (moyennant le yoga général des motifs --- il est
possible qu'on utilise pour ceci les conjectures de Hodge …) que le groupe
$H'$ est une limite projective de groupes algébriques \struck{\ill{}}
unipotents \ill{}. Si on se borne à une catégorie de motifs à engendrement
fini, et si $G$ est connexe, on trouve aussi que $H'$ est connexe, donc étant
résoluble, on trouve}
\note{Tout ce N.B.\ est barré d'une dizaine de longues diagonales ; il se lit
encore, et le paragraphe qui suit le remplace.}

On constate aisément que le groupe $H'$ est nécessairement une limite
projective de groupes algébriques unipotents. On en conclut aussitôt que si
$\underline{C}$ est à engendrement fini (ou, plus généralement, \add{à
engendrement} dénombrable, de façon que $G'$ donc aussi $H'$ soit limite
projective d'une suite de groupes algébriques) alors tout torseur sous $H'$
est trivial ; cela signifie ici que tout foncteur fibre filtré est en fait
associé à un foncteur fibre gradué (en prenant la filtration correspondant à
la graduation),

\page{105}
i.e.\ que la filtration dudit foncteur admet un splittage compatible avec les
produits tensoriels. Mais le choix d'un tel splittage équivaut à celui d'un
point sur un certain torseur à droite $Q = Q'$ sous le \add{schéma en}
groupe \add{(pro-unipotent)} $H$ des automorphismes de $F$ respectant la
filtration et induisant l'identité sur le gradué associé ; il n'est pas du
tout canonique !

3) Appelons \emph{pré-structure de Hodge} sur un espace vectoriel \add{$V$} de
dimension finie sur $\underline{\mathbb{Q}}$, la donnée d'une bigraduation sur
$V \otimes_{\underline{\mathbb{Q}}} \underline{C} = V_{\underline{C}}$,
\[
V_{\underline{C}} = \coprod_{p,q} V^{p,q} ,
\]
telle que a) la graduation totale correspondante soit « définie sur
$\underline{\mathbb{Q}}$ » i.e.\ $\coprod_{p+q=n} V^{p,q}$ provienne d'un
sous-espace $V_{\underline{\mathbb{Q}}}^{n}$ de $V$, et b) on a
$\overline{V^{p,q}} = V^{q,p}$, où $x \mapsto \bar{x}$ désigne la
conjuguaison complexe. (N.B.\ Généralisation à des corps plus généraux laissée
à Saavedra). Les vectoriels $V$ munis d'une pré-structure de Hodge forment une
catégorie tensorielle sur $\underline{\mathbb{Q}}$ dans un sens évident, muni
d'un foncteur fibre canonique, le foncteur « oubli » \add{$F$}. On trouve donc
un groupe de Galois $G$, et plus généralement toute \add{$\otimes$-}sous-catégorie
$\underline{C}$ de la catégorie précédente nous définit un groupe de Galois
$G$. La bigraduation sur le foncteur
$F_{\underline{C}}(V) = V \otimes_{\underline{\mathbb{Q}}} \underline{C}$
correspond, en vertu de 1) (où il convient cependant de se permettre une
extension du corps de base sur le foncteur fibre envisagé …) d'un
homomorphisme
$\underline{G}_{m\,\underline{C}}^{2} \to G_{\underline{C}}$, i.e.\ de deux
homomorphismes \add{qui commutent}
\[
i_{1}, i_{2} : \underline{G}_{m\,\underline{C}} \longrightarrow G_{\underline{C}} .
\]
Les deux conditions a) et b) imposées aux structures envisagées
s'interprètent respectivement par les faits que l'homomorphisme
\[
i_{\underline{C}} = i_{1} i_{2} : \lambda \mapsto i_{1}(\lambda) i_{2}(\lambda)
\]
est « défini sur $\underline{\mathbb{Q}}$ » i.e.\ provient d'un homomorphisme
\[
i : \underline{G}_{m} \longrightarrow G ,
\]
(nécessairement central, car les composantes homogènes
$V_{\underline{\mathbb{Q}}}^{n}$ d'une pré-structure de Hodge sont
\struck{nécessairement} \add{évidemment} munis d'une \add{pré-}structure de
Hodge de façon que $V$ soit la somme directe des
$V_{\underline{\mathbb{Q}}}^{n}$ en tant que pré-structure

\page{106}
de Hodge), et par la condition que l'on a
\[
i_{2} = \overline{i_{1}} \quad \text{i.e.} \quad
i_{2}(\lambda) = \overline{i_{1}(\bar{\lambda})} \quad \text{pour tout}\quad
\lambda \in \underline{C} .
\]

Si à toute variété projective lisse $X$ sur $\underline{C}$ on associe sa
cohomologie rationelle $H^{*}(X, \underline{\mathbb{Q}}) = V$, de sorte que
$V_{\underline{C}} = H^{*}(X, \underline{C})$ est isomorphe canoniquement (par
la théorie de Hodge) à
$H_{Hdg}(X) = \coprod H^{q}(X, \underline{\Omega}^{p}_{X/\underline{C}})$, on
voit qu'on trouve ainsi une pré-structure de Hodge sur
$H^{*}(X, \underline{\mathbb{Q}})$, d'où un \add{$\otimes$-}foncteur naturel
de la catégorie des motifs sur $\underline{C}$ dans la catégorie des
pré-structures de Hodge. La conjecture de Hodge standart équivaut à dire que
ce foncteur est \emph{pleinement fidèle}, i.e.\ que l'homomorphisme naturel
qui va du groupe de Galois de Hodge précédent $G$ dans le groupe de Galois
motivique (associé au \add{$\otimes$-}foncteur de Betti $H_{Bet}$) est un
épimorphisme. Ou encore que pour toute catégorie \add{$\underline{C}_{0}$} de
motifs \add{de type fini} sur $\underline{C}$, désignant par $\underline{C}$
la \add{$\otimes$-}catégorie de pré-structures de Hodge engendrée par les
$H_{Bet}(M)$ pour $M \in \operatorname{Ob}(\underline{C}_{0})$,
l'homomorphisme de groupes algébriques $G \to G_{0}$ associé au foncteur de
Betti-Hodge $\underline{C}_{0} \to \underline{C}$ est un épimorphisme (i.e.\
surjectif sur les points à valeurs complexes, disons).

On appelle \struck{\ill{}} \emph{polarisation} d'une pré-structure de Hodge de
poids $n$ la donnée d'un accouplement de préstructures de Hodge
\[
\varphi : V \times V \longrightarrow \underline{\mathbb{Q}}(n) ,
\]
où $\underline{\mathbb{Q}}(n)$ est l'espace vectoriel trivial
$\underline{\mathbb{Q}}$ sur $\underline{\mathbb{Q}}$, avec
$\underline{\mathbb{Q}}_{\underline{C}} = \underline{C}$ muni du bidegré
$(n,n)$), ayant la propriété que la forme hermitienne correspondante sur
$V_{\underline{C}}$,
\[
\psi(x,y) = \varphi(x, \bar{y})\, (-i)^{p-q}
\quad \text{pour $x$ de bidegré } (p, n-p)
\]
\note{Le bidegré est tapé $(p,q)$ ; le $n$ est écrit à l'encre par-dessus le
$q$.}
soit définie positive. Une \emph{structure de Hodge} est une préstructure de
Hodge admettant une polarisation. Les structures de Hodge forment une
sous-$\otimes$-catégorie de la catégorie des pré-structures de Hodge. La
théorie de Hodge nous assure que le foncteur de Betti-Hodge sur la catégorie
des

\page{107}
motifs sur $\underline{C}$ prend ses valeurs en fait dans la catégorie des
structures de Hodge (une polarisation d'une \uncertain{varéiét}\note{sic}
projective lisse $V$ définit canoniquement une polarisation de la structure de
Hodge associée sur la cohomologie de Betti-Hodge). N.B.\ On n'a aucune idée
sur ce que pourrait être l'image essentielle du foncteur précédent, par
exemple s'il y a lieu d'espérer qu'on trouve toutes les structures de Hodge
(donc une équivalence de catégories : motifs sur $\underline{C} \to$
structures de Hodge) ; cela semble peu probable, mais on n'a aucune indication
sérieuse dans un sens ou l'autre.

La catégorie des pré-structures de Hodge est semi-simple. Si $G$ est le groupe
de Galois d'une \add{$\otimes$-}catégorie \add{$\underline{C}$} de
pré-structures de Hodge, à engendrement fini si on veut (pour simplifier),
alors on peut expliciter en termes du groupe de Galois associé et de sa
structure $i_{1}$ ci-dessus la condition pour que les objets de
$\underline{C}$ soient en fait des structures de Hodge. Ceci est un exercice
plaisant et délectable, qui devrait figurer dans un travail systématique sur
les $\otimes$-catégories, dans le chapitre des exemples. On trouve des
restrictions très sérieuses sur le groupe $G$ muni de $i_{1}$ (en plus du fait
que $G$ soit réductif).

4) Je laisse le soin à Saavedra de déterminer quelle structure supplémentaire
on obtient sur la structure de Hodge « complexe » associée à une variété
projective \add{lisse} complexe $X$, lorsqu'on se donne cette dernière comme
déduite d'une variété projective réelle $X_{\underline{R}} = X_{0}$. On trouve
une notion de « structure de Hodge réelle », donnant naissance à une
$\otimes$-catégorie correspondante. Dans le groupe de Galois motivique de
celui-ci, en plus de la structure $i_{1}$, on trouve un élément $f_{\infty}$
de $G(\underline{\mathbb{Q}})$, \struck{de degré} d'ordre $2$ (jouant le rôle
d'un « élément de Frobenius à l'infini »), qui correspond à l'automorphisme du
foncteur de Betti
$X_{0} \longmapsto H^{*}(X_{0}(\underline{C}), \underline{\mathbb{Q}})$ déduit
de l'homéomorphisme $x \mapsto \bar{x}$ de $X_{0}(\underline{C})$. Il faut
expliciter les relations entre cet élément et $i_{1}$, $i_{2}$ !

\page{108}
5) : Soit $\underline{C}$ une $\otimes$-catégorie tensorielle,
\struck{\ill{}} munie d'un foncteur fibre sur $k$ \add{de caractéristique
nulle.} A prouver que, pour que le groupe de Galois $G$ correspondant soit
profini, il faut et il suffit que pour tout objet $M$ de $\underline{C}$, la
$\otimes$-catégorie engendrée soit semi-simple et n'ait qu'un nombre fini
d'objects simples non isomorphes. Si $\underline{C}$ est quelconque, la
sous-catégorie pleine $\underline{C}_{0}$ de $\underline{C}$ qui correspond au
pro-groupe quotient de $G$ formé des $G_{i}^{\circ}$ (où
$G = \varprojlim G_{i}$, et $G_{i}^{\circ}$ est la composante neutre de
$G_{i}$) est formée exactement des objets $M$ ayant la propriété précédente.

Il serait intéressant de trouver des énoncés correspondants en caractéristique
quelconque.

6) La notion de polarisation d'un motif sur un corps (elle-même déduite de
celle de polarisation d'une variété algébrique) donne une structure
supplémentaire remarquable dans la catégorie des motifs : si $M$ est un motif
de poids $n$, on sait parmi les formes symétriques ($n$ pair) resp.\ alternées
($n$ impair) $M \otimes M \to T(n)$ (où $T$ est le motif de Tate) distinguer
celles qui sont « définies positives » ou encore des « polarisations ». Cette
notion se reflête par exemple par des structures supplémentaires sur les
groupes de Galois motiviques. Il y a lieu de faire une étude axiomatique
abstraite d'une telle notion de polarisation sur une $\otimes$-catégorie
générale au dessus d'un sous-corps du corps des réels. On pourra en rediscuter
à l'occasion.

\section*{Catégories de modules sur un schéma en groupes affine, de sa main (pages 110 à 112)}

\note{Trois feuillets entièrement de sa main, écrits vite. Les formules
passent ; la prose qui les relie ne passe pas toujours, et ce qu'elle ne
donne pas est marqué plutôt que comblé.}

\page{110}
\struck{$G$ schéma en \ill{} sur le corps $k$} \add{un gr.\ \ill{}, qcq,
quasi-sép.\ et plat sur une base}
\note{La première ligne est reprise en interligne : la restriction au corps
$k$ est barrée et remplacée par des hypothèses relatives sur une base. Le
premier mot de l'insertion ne se laisse pas lire.}

Tout vectoriel $V$ sur $k$ où $G$ opère est \ill{} des sous-vectoriels
$V_{i}$ de dim.\ finie sur $k$ \ill{}.

$G$ schéma en groupes sur $S$. \ill{} conditions :
\begin{enumerate}
\item[a)] $G$ affine ;
\item[b)] $G$ \uncertain{qu.-cpt.\ quasi-sép.} sur $S$ et plat sur $S$.
\end{enumerate}

\marginal{N.B. Il suffit que $A(G)$ soit \ill{} pour « loc » \ill{}}
\note{Cette note, cerclée dans la marge gauche, ne se lit que par bribes ; on
ne complète pas.}

Soit $\underline{C}_{qc}(G)$ la catégorie des $S$-modules qu.\ coh.\ où $G$
opère, $\underline{C}_{f}(G)$ la sous-catégorie pleine formée des
$\underline{E}$ qui sont de t.f. Alors
\[
\underline{C}(G) \overset{\approx}{\longrightarrow} \operatorname{Ind} \underline{C}_{f}(G)
\]
cf.\ SGAD VI$_{B}$ 11.2.

Soit alors $f : T \to S$ un morphisme affine. \add{Nous poserons}
$\mathcal{B} = f_{*}(\mathcal{O}_{T})$, \struck{\ill{}} défini par
\struck{$A$-modules} définir la catégorie des $G_{T}$-\struck{$E$-}Modules
quasi-cohérents en termes de $\underline{C}(G)$ et $\mathcal{B}$.

Or la donnée d'un $\mathcal{O}_{T}$-module $\underline{F}$ qu.\ coh.\ revient
à la donnée d'un $\underline{E}$ (qu.\ coh.\ sur $S$,
$\bigl(= f_{*}(\underline{F})\bigr)$) sur lequel $\mathcal{B}$ opère. On
constate de plus (avec hyp.\ sur $G$) que la donnée d'une opération

\page{111}
de $G_{T}$ sur $\underline{F}$ équivaut à la donnée d'une opération de $G$ sur
$\underline{E}$ \add{commutant à $\mathcal{B}$}, car
$G_{T} \to \underline{\operatorname{Aut}}(\underline{F})$ équivaut à
\[
G \longrightarrow \prod_{T/S} \underline{\operatorname{Aut}}(\underline{F}) ,
\quad \text{or}
\]
\[
\prod_{T/S} \underline{\operatorname{Aut}}(\underline{F}) \simeq
\underline{\operatorname{Aut}}_{\mathcal{B}}\Bigl(\prod_{T/S} \underline{F}\Bigr)
= \underline{\operatorname{Aut}}_{\mathcal{B}}(\underline{E}) .
\]

Donc, la catégorie $\underline{C}(G_{T})$ équivaut à la catégorie des objets
de $\underline{C}(G)$ munis d'une opération de $\mathcal{B}$, ce qui, si $S$
est affine d'anneau $A$, donc $\mathcal{B}$ défini par $A$-algèbre $B$,
équivaut à la \struck{catég.} donnée d'un objet de $\underline{C}(G)$ à anneau
$B$ d'endomorphismes. Lorsque $G$ satisfait l'une des conditions a) ou b)
(p.~ex.\ $G$ \ill{} sur un corps $k$), la catégorie $\underline{C}(G_{T})$
est définie en termes de $\underline{C}_{f}(G)$. De même, si $\varphi_{G}$
désigne le foncteur oubli de $\underline{C}_{f}(G)$ \add{$\varphi_{G}^{f}$}
dans \ill{},
$\bigl(\underline{C}(G_{T}), \varphi_{G_{T}}\bigr)$ est déterminé en termes de
$\bigl(\underline{C}_{f}(G), \varphi_{G}^{f}\bigr)$. Et il en est de même de
$\bigl(\underline{C}_{f}(G_{T}), \varphi_{G_{T}}\bigr)$, sous-catégorie pleine
de $\underline{C}(G_{T})$ formée des $E$ tels que $\varphi_{G_{T}}(E)$
\uncertain{de d.f.}\note{ainsi écrit ; la page 110 porte « t.f. » au même
emploi} sur $\mathcal{B}$.

\page{112}
\marginal{\uncertain{loc. lis.}}

Soit $G$ un schéma \emph{affine} sur un anneau $A$. Dans quelle mesure $G$ se
récupère-t-il par la connaissance de
$\bigl(\underline{C}_{f}(G), \varphi_{G}^{f}\bigr)$, i.e.\
$\underline{C}_{f}(G)$ est regardé comme une $\otimes$-catégorie
\add{par $\otimes$ \ill{}}, et $\varphi_{G}^{f}$ comme un
$\otimes$-foncteur.

Il faudrait récupérer $G(B)$ pour les $A$-algèbres $B$. Or \ill{}, il
suffirait de récupérer $G(A)$. On a un hom.\ canonique
\[
G(A) \longrightarrow \operatorname{Aut}_{\otimes}(\varphi_{G}^{f})
\]
\struck{\ill{}}

On prouve que c'est injectif : regardons les effets de $g \in G(A)$ sur
l'\uncertain{algèbre} \add{affine $E$} de $G$ !

On prouve sans doute aussi que c'est surjectif, en regardant comme l'effet
d'un $\otimes$-\ill{} \ill{} de $\varphi_{G}^{f}$ sur $E$,
\struck{\ill{}} \marginal{à détailler}

\struck{\ill{}}

\emph{Question} : À quelle condition une $\otimes$-catégorie
\add{$A$-linéaire} munie d'un $\otimes$-foncteur
$\varphi^{f} : \underline{C}_{f} \to \operatorname{Mod}_{f}(A)$ provient-elle
d'un schéma en groupes affine sur $A$ ???

Il faudrait supposer que $\underline{C}$ est un $\mathcal{A}b$-topos, et que
$\varphi$ commute aux $\varinjlim$.

\section*{Plan du traité (page 114)}

\note{Feuillet tapé, complété de sa main : l'appendice du chapitre 1 et les
deux derniers alinéas sont à l'encre.}

\page{114}
\begin{itemize}
\item Chapitre 1. Généralités sur les opérations $\otimes$ sur une catégorie.
  \add{Appendice : Rappels sur les coalgèbres, comodules, bigèbres …}
\item Chapitre 2. Relations entre $\otimes$-catégories \struck{et} avec
  foncteur fibre, et coalgèbres resp.\ hyperalgèbres : \struck{Groupe} Dualité
  de Tannaka. (structures mixtes, connexions …)
\item Chapitre 3. Exemples nombreux et divers. Dictionnaire détaillé.
\item Chapitre 4. Filtrations sur un foncteur fibre.
\item Chapitre 5. Théorème de structure des $\otimes$-catégories en termes de
  gerbes.
\item Chapitre 6. Polarisation d'une $\otimes$-catégorie. Groupes de Hodge.
\item Chapitre 7. $\otimes$-Catégories à lien réductif. Classification sur
  $\underline{R}$.
\item Chapitre 8. Structures de Hodge, groupes de Mumford-Tate et espaces
  modulaires correspondants.
\item Chapitre 9. Questions d'effectivité (?)
\end{itemize}

\marginal{Éventuellement un chapitre de structure sur le « cas non commutatif »}

\marginal{Un autre, une variante pour catégories triangulées ?}

\marginal{Chapitre sur catégories de Picard, (plus gén.\ gpoïdes avec loi
$\otimes$) catégories de Picard enveloppantes. Cas des catégories ayant deux
lois $\oplus$ et $\otimes$ ----}

\section*{Table des matières du chapitre 1 (pages 115 à 116)}

\note{Table tapée, annotée partout : les numéros de paragraphes sont
renumérotés à l'encre --- le Par.\ 3 devient 4, le Par.\ 4 devient 5, et ainsi
de suite --- deux subdivisions sont ajoutées de sa main, une est barrée. On
donne la numérotation corrigée.}

\page{115}
\subsection*{Chapitre 1. Généralités sur les $\otimes$-catégories.}

\marginal{\struck{\ill{}} --- 1) produit tens.\ de modules ; 2) catégorie
avec produits \uncertain{cart.} ; 3) Torseurs sous un faisceau abélien ;
4) familles de modules \ill{} indexées par \ill{} ; 5) espace des
\uncertain{lacets} d'un \ill{} \add{topologique}}

Par.~0 : Lois $\otimes$ sur une catégorie : exemples divers …

Par.~1 \struck{I. Les données} Contraintes pour une loi $\otimes$.
\begin{itemize}
\item 1.1. Associativité.
\item 1.2. Commutativité.
\item 1.3. Unités.
\item 1.4. $\underline{\operatorname{Hom}}(X,Y)$,
  $\underline{\operatorname{Moh}}(Y,X)$, \struck{\ill{}}.
  \note{« Moh » est bien ce que porte la machine, et le mot revient plusieurs
  fois : c'est le Hom interne pris de l'autre côté.}
\item 1.5. \add{$\check{X}$, \struck{$\check{X}^{-1}$} $X^{-1}$.}
\item 1.6. \add{Cas des exemples du Par.~0}
\end{itemize}

Par.~2 Compatibilités entre contraintes.
\marginal{Toutes les compatibilités entre contraintes sont satisfaites, quand
elles ont un sens, dans les exemples du par.~0}
\begin{itemize}
\item 1.1. Associativité et commutativité.
\item 1.2. Commutativité et unité.
\item 1.3. Associativité et unité.
\item 1.4. $\underline{\operatorname{Hom}}$,
  $\underline{\operatorname{Moh}}$ et associativité : formules de Cartan.
\item 1.5. $\underline{\operatorname{Hom}}$ et commutativité (on a alors
  $\underline{\operatorname{Hom}}(X,Y) \simeq
  \underline{\operatorname{Moh}}(Y,X)$) … \marginal{non !}
\item 1.6. \struck{$\underline{\operatorname{Hom}}$ et unité} : définition de
  $\check{X}$ (pour loi commutative seulement ?). \marginal{Lois $\otimes$
  rigides.}
\item 1.7. Formalisme des traces (cas d'isomorphie
  $\check{X} \otimes Y \simeq \underline{\operatorname{Hom}}(X,Y)$),
  \marginal{rang d'un objet.}
\end{itemize}

Par.~3 Formes normales de multifoncteurs composés. \struck{\ill{}}

Par.~4 $\otimes$-foncteurs.
\begin{itemize}
\item 4.1. Théorie sans contraintes : définition des $\otimes$-foncteurs,
  composition, morphismes de $\otimes$-foncteurs.
\item 4.2. Compatibilité \struck{d'un} $\otimes$-foncteur avec des
  contraintes : associativité et commutativité.
\item 4.3. Compatibilité avec des contraintes : unité.
\item 4.4. Compatibilité avec des contraintes :
  $\underline{\operatorname{Hom}}$, $\underline{\operatorname{Moh}}$,
  $\check{X}$, $X^{-1}$.
\item 4.5. \add{Objets quasi-inversibles, et \struck{\ill{}} calcul des
  fractions …}
\end{itemize}

Par.~5 Exemples typiques ; \struck{les} $\otimes$-catégories \add{définies par
des bigèbres.}
\begin{itemize}
\item 5.0. \add{$\otimes$-Catégories : catégories additives, avec $\otimes$
  biadditif, ass.\ comm.\ unit.\ \struck{\ill{}} (pas néc.\
  $\underline{\operatorname{Hom}}$).}
\item 5.1. Modules sur une bigèbre.
\item 5.2. Comodules sur une bigèbre.
\item 5.3. Dualité entre les deux notions : première manière (imparfaite).
\item 5.4. Dualité entre les deux notions : seconde manière, formelle (via
  \struck{\ill{}} notions d'algèbre dans une $\otimes$-catégorie).
\item 5.5. \add{Modules mixtes (co $+$ ordinaire) sur une bigèbre}
\end{itemize}

\page{116}
\begin{itemize}
\item 5.6. Variantes avec bi-proalgèbres.
\item 5.7. Variantes graduées.
\end{itemize}

Par.~6. Bigèbres d'endomorphismes des $\otimes$-foncteurs. Schémas en groupes
d'automorphismes de $\otimes$-foncteurs.

\ill{} …

Par.~7. Structure des catégories de Picard \add{(?)} (Viendrait mieux avant
par.~4). \marginal{N'est pas nécessairement le \uncertain{lieu} ici}

\marginal{Par.\ \struck{\ill{}} Plus gén., \struck{catégories} monoïdes avec
$\otimes$ associatif \uncertain{top.\ an.} et unitaire (pas néc.\ comm.\
…).}

\section*{$\otimes$-catégories (pages 117 à 120)}

\note{Quatrième tapuscrit du lot, corrigé de sa main. La machine écrit $\phi$
par un o barré et $\psi$ par un y barré ; on donne les lettres grecques.}

\page{117}
\subsection*{$\otimes$-catégories.}

1) \emph{Définition}. Catégorie \struck{\ill{}} additive, avec bifoncteur
\[
(X,Y) \mapsto X \otimes Y
\]
\struck{et} un \struck{objet} isomorphisme de bifoncteurs
\[
\varphi(X,Y,Z) : (X \otimes Y) \otimes Z \overset{\sim}{\longrightarrow}
X \otimes (Y \otimes Z)
\]
et un isomorphisme de commutativité
\[
\psi(X,Y) : X \otimes Y \overset{\sim}{\longrightarrow} Y \otimes X
\]
satisfaisant des axiomes
\begin{enumerate}
\item[a)] Axiome du pentagone pour $\varphi$ \struck{(ne dépend pas de la
structure additive)}
\item[b)] $\psi(X,Y)\,\psi(Y,X) = \operatorname{id}_{Y \otimes X}$ ;
\item[c)] compatibilité entre $\varphi$ et $\psi$ : \add{axiome de
l'hexagone}
\item[d)] existence d'un élément $\underline{1}_{\underline{C}} = \underline{1}$,
avec des isomorphismes fonctoriels
\add{$v(X) = $}
\[
u(X) : \underline{1} \otimes X \simeq X \qquad \text{d'où} \quad
\psi(X,\underline{1})\,u(X) : X \otimes \underline{1} \simeq X ,
\]
\struck{(donc par b) on aura aussi $u(X) = v(X)\,\psi(\underline{1},X)$)}
avec la condition que pour $X = \underline{1}$, on trouve le même
isomorphisme
\[
u(\underline{1}) = v(\underline{1}) : \underline{1} \otimes \underline{1}
\simeq \underline{1}
\]
ce qui s'exprime aussi par la condition que
$\psi(\underline{1},\underline{1}) = \operatorname{id}_{\underline{1}}$ ;
\marginal{(de plus on veut) compatibilité de $(\underline{1},u,v)$ avec
$\varphi$, quand $X, Y$ \ill{} $2$ \ill{} i.e.\ $2 = 1$, avec $\psi$ \ill{}
$X$ --- $Y$ est $= \underline{1}$.}
\item[e)] le foncteur $\otimes$ est biadditif.
\end{enumerate}

\marginal{Envisager aussi la structure éventuelle
$\underline{\operatorname{Hom}}(X,Y)$, relation avec $\check{X} \otimes Y$}

N.B. a) \struck{et d) sont} \add{est} indépendants de la donnée de $\psi$,
b) est indépendant de la donné de $\varphi$, a) b) c) d) sont indépendant du
fait que $\underline{C}$ est additive.

Soit $\underline{1}', u', $\add{$v'$} comme $\underline{1}, u, v$, on
\struck{va définir canoniquement un} \add{voit qu'il existe un unique}
isomorphisme $w : \underline{1} \simeq \underline{1}'$, compatible en un sens
évident avec $u, v$ et $u', v'$. On choisit le composé
\[
\underline{1} \xrightarrow{\ v'(\underline{1})\ } \underline{1} \otimes
\underline{1}' \xrightarrow{\ u(\underline{1}')\ } \underline{1}'
\]
\note{Les deux étiquettes sont ajoutées à la main au-dessus des flèches ; la
fin de la ligne est barrée d'une dizaine de traits.}
A vérifier les compatibilités annoncées (ne dépend pas de la donnée de
$\varphi$, \add{ni $\psi$,} mais de $u, v$ satisfaisant la compatibilité
$u(\underline{1}) = v(\underline{1})$ et $u', v'$ kif-kif). C'est pourquoi il
n'y \add{a} pas lieu de poser $\underline{1}, u, v$ comme donnée, mais comme
axiome d'existence.

\page{118}
2) $\otimes$-Foncteur $F : \underline{C} \to \underline{C}'$ (avec
$\underline{C}, \underline{C}'$ des trucs-catégories).
\note{« trucs-catégories » est bien ce que porte la machine ; le mot tient
lieu de « $\otimes$-catégories ».}

Les données sont un foncteur \struck{\ill{}} $F$, plus \struck{un}
isomorphismes de bifoncteurs
\[
c(X,Y) : F(X \otimes Y) \overset{\sim}{\longrightarrow} F(X) \otimes F(Y)
\]
\struck{\ill{}} satisfaisant les conditions :
\begin{enumerate}
\item[a)] Compatibilité \add{de $c$} avec la donnée de $\varphi$, $\varphi'$
\item[b)] Compatibilité \add{de $c$} avec la donnée de $\psi$, $\psi'$

Existe isom.\ \struck{foncto.} $i : F(\underline{1}) \simeq \underline{1}'$
compatible avec $c$ et $u'$
\item[c)] \struck{Compatibilité de $i$ avec \ill{}} avec les isomorphismes
$u$ (ou ce qui revient au même, avec \add{$c$ et} $v, v'$).
\end{enumerate}

N.B. Cet $i$ sera alors uniquement déterminé.

\begin{enumerate}
\item[d)] $F$ est additif.
\end{enumerate}

3) Homomorphisme d'un $\otimes$-foncteur $F$ dans un autre $G$ : c'est un
homomorphisme de foncteurs $f : F \to G$ satisfaisant la condition
\begin{enumerate}
\item[a)] compatibilité avec \struck{\ill{}} $c_{F}$, $c_{G}$
  \marginal{\struck{\ill{}}} \marginal{sans doute pas une}
\item[b)] \struck{b) --- (i)} compatibilité avec $i_{F}$, $i_{G}$
  (\add{ça} c'est \struck{peut-être} conséquence formelle de a) --- à
  regarder), \add{cf.\ plus bas)}
\end{enumerate}

4) \emph{Remarques}. A dissocier \add{ce qui précède} en plusieurs aspects de
la donnée d'un bifoncteur $\otimes$ :
\begin{enumerate}
\item[a)] Donnée d'associativité, soumise à l'axiome du pentagone. Un foncteur
« compatible avec $\otimes$ \add{associatif} » de $\underline{C}$ dans
$\underline{C}'$ est un foncteur $+$ un $c$, compatible avec la donnée de
$\varphi$.
\item[b)] Donnée d'unité bilatère $\underline{1}, u, v$ (unique à isomorphisme
unique près). Un foncteur « compatible avec \add{$\otimes$ à} unité bilatère »
est un $F$ muni de $c$, tel qu'il existe $i$ compatible avec $c, u, v$ et
$c', u', v'$ (lequel $i$ sera alors unique).
\marginal{N.B. Le monoïde des endom.\ de $\underline{1}$ est alors
commutatif.}
\item[c)] Donnée de commutativité $\psi$, satisfaisant la condition de
réflexivité. Un foncteur
\end{enumerate}

\page{119}
« compatible avec $\otimes$ commutatif » est un foncteur et un $c$, compatible
avec la donnée des $\psi$.

Quand on a une loi $\otimes$ avec donnée a) et b), il y a lieu de postuler une
compatibilité entre les deux (dans le cadre de l'axiomatique 1), est-ce une
conséquence des autres axiomes ? je \struck{crains} \add{pense} que non, et
qu'il faille le rajouter). Quand on a une loi $\otimes$ avec donnée a) et
\uncertain{c}), il y a lieu de postuler une compatibilité entre les deux. De
même pour une loi $\otimes$ avec b) et c). Quand on a à la fois a) b) c), le
paquet \struck{\ill{}} de toutes les conditions de compatibilité sont celles
énoncées plus haut \add{dans 1°} a) à d). De même, un Foncteur
« compatible avec $\otimes$ associatif à unité bilatère » est un foncteur $F$,
$+$ un $c$, compatible séparément avec associativité et unité bilatèrité, et
de même pour tout autre paquet de données supplémentaires parmi a), b), c).
Enfin, la notion d'homomorphisme de \struck{\ill{}} $\otimes$-foncteurs ne
dépend pas des structures supplémentaires a) \struck{b)} ou c)
\add{(mais bien de b))} qu'on peut envisager, s'exprimant exclusivement à
l'aide de la donnée de c) \struck{(sous réserve de vérification indiquée dans
3°)}. --- Quant à la \struck{\ill{}} \add{biadditiv\ill{}}
condition d'additivité sur les catégories et de biadditivité sur $\otimes$, il
n'y a lieu de l'introduire qu'en dernier lieu.

\marginal{Attention si l'associativité ne doit pas rappliquer ici aussi}

5) N.B. Si on a $\otimes$ avec unité bilatère, alors pour deux endomorphismes
$u, v$ de $\underline{1}$, on a $u \otimes v = uv = vu$. Donc sous la
condition de commutativité c), $\operatorname{End}(\underline{1})$ est un
monoïde commutatif. \struck{\ill{}} Sans condition de comm.\ ce monoide opère
à gauche et à droite sur tout objet de $\underline{C}$ ; les deux opérations
coïncident moyennant commutativité. Si \struck{la structure}
\add{$\underline{C}$} est additive, $\operatorname{End}(\underline{1})$ est un
anneau, qui opère à gauche et à droite sur chaque objet de $\underline{C}$.
Dans le cas commutatif, $\operatorname{End}(\underline{1}) = k$ est donc un
anneau commutatif et la catégorie $\underline{C}$ est $k$-linéaire de façon
naturelle, \add{$\otimes$ est $k$-bilinéaire.}

\page{120}
6) Supposons maintenant que $\underline{C}$ soit une catégorie
\struck{additive} \add{abélienne} $k$-linéaire ($k$ un corps) munie d'un
foncteur fidèle exact $F : \underline{C} \to \operatorname{Modf}(k)$, de sorte
que $\underline{C}$ s'identifie à la catégorie des représentations
\struck{\ill{}} \add{(continues)} finies d'une algèbre profinie $U$ sur $k$
\add{(associative, unitaire)}, ou des comodules sur le dual
\add{topologique} $A$ de $k$, qui est une coalgèbre (associative, unitaire).
Alors la donnée d'un foncteur $\otimes$ \add{bilinéaire} compatible avec $F$
(i.e.\ pour lequel \struck{\ill{}} on a \emph{égalité}
$F(X \otimes Y) = F(X) \otimes F(Y)$) \emph{équivaut} à la donnée d'un
morphisme diagonal $U \to U \mathbin{\widehat{\otimes}_{k}} U$, ou encore
d'une loi de multiplication sur $A$.

\marginal{généraliser au moins partiellement au cas d'un $k$ qui n'est pas un
corps}
\note{Cette phrase est tapée dans la marge gauche, en regard des lignes qui
précèdent.}

Une donnée supplémentaire (au sens de 4)) a) \struck{b)} ou c)
\emph{compatible} avec $F$, $c = \operatorname{id}$ est alors uniquement
déterminée, si elle existe.
\struck{pour b) c'est une \ill{} $e : U \to k$ (morph.\ d'alg.) i.e.\
$\varepsilon \in A$ \ill{}}
Elle existe sss $A$ est associative, resp.\ unitaire, resp.\ commutative.
\note{« sss » est la graphie de la machine pour « si et seulement si » ; elle
revient plusieurs fois dans ce lot.}
Quand ces conditions sont satisfaites simultanément, on a donc un schéma en
\struck{groupes} \add{monoïdes} affine sur $k$, qui permet de réconstituer la
situation. --- Si on a de même $\underline{C}', F'$ sur le même $k$,
\struck{\ill{}} alors les foncteurs $G : \underline{C} \to \underline{C}'$
compatibles avec $F, F'$ correspondent exactement aux hom.\ de $k$-algèbre top
$U' \to U$, ou de $k$-coalgèbres $A \to A'$. Si $\underline{C}$ et
$\underline{C}'$ sont munis de structures bilinéaires $\otimes$, alors le
foncteur $G$ est un $\otimes$-foncteur avec un isomorphisme $c$ transformé par
$F'$ en $\operatorname{id}$ (ce qui suffit à \struck{le} caractériser
\add{alors}) sss $U' \to U$ est compatible avec les diagonales i.e.\
$A \to A'$ est compatible avec les multiplications. Les éventuelles structures
supplémentaires a), \struck{b)} ou c) sont respectées automatiquement,
\marginal{pour b) il faut supposer compatible avec augment.\ de $U$ resp.\
avec unité de $A$, $A'$}
\marginal{Dire que les endom.\ de $F$ compatibles avec
$c = \operatorname{id}$ sont les $g \in U$ tels que
$\Delta(g) = g \otimes g$ ; si on veut respecter la donnée de l'unité de
$(\underline{C}, \otimes)$, il faut \ill{} la cond.\ $\varepsilon(g) = 1$,
i.e.\ \ill{} les morphismes $A \to k$}
\marginal{relations entre loi $\check{X}$ et antipo\add{de} données par $U$
resp.\ $A$}

7) \emph{Un exemple}. On prend pour $\underline{C}$ la catégorie des modules
gradués sur un anneau commutatif $k$, avec le produit tensoriel gradué
ordinaire et pour unité le module $k_{s}$ gradué par le degré zéro, données de
$u$ et $v$ comme d'habitude, ainsi que $\varphi$ (associativité) et $\psi$
(commutativité). Tous les axiomes du n° 1 sont vérifiés, et le
foncteur oubli (avec $c = \operatorname{id}$) est un $\otimes$-foncteur
respectant toutes les données a) b) c). Donc si $k$ est
\note{La phrase s'arrête là : le feuillet est le dernier du lot.}

\end{document}
